% main.tex — Politecnico di Torino Beamer presentation template demo/guide
\documentclass{beamer}

\usepackage{theme/beamerthemesintef} % PoliTo theme
\usepackage{amsfonts,amsmath,oldgerm}

\usetheme{sintef}
\AtBeginSection[]{} % Disable auto-generated TOC slides

% ── Enhanced visual packages ──────────────────────────────────────────────────
\usepackage{fontawesome5}       % Vector icons (GitHub, tools, arrows, etc.)
\usepackage[most]{tcolorbox}    % Custom styled boxes
\usepackage{booktabs}           % Professional tables
\usepackage{tabularx}           % Flexible-width table columns
\usepackage{tikz}               % Native diagrams
\usetikzlibrary{arrows.meta, positioning, shapes.geometric, calc, fit, backgrounds}
\tikzset{
    layer/.style={draw=#1, fill=#1!8, rounded corners=2pt,
        minimum width=5.5cm, minimum height=0.55cm,
        font=\scriptsize, align=left, text=darkgray}
} % Defined here (not inside frames) so #1 is handled by pgfkeys, not Beamer
\usepackage{listings}           % Code listings
\usepackage{xcolor}             % Extended color support
\usepackage{hyperref}           % Clickable links
\usepackage{multicol}           % Multi-column layouts

% ── Theme helpers ─────────────────────────────────────────────────────────────
\newcommand{\testcolor}[1]{\colorbox{#1}{\textcolor{#1}{test}}~\texttt{#1}}
\usefonttheme[onlymath]{serif}
\titlebackground*{theme/assets/background}
\newcommand{\hrefcol}[2]{\textcolor{cyan}{\href{#1}{#2}}}

% ── Custom colors (Demo Showcase) ─────────────────────────────────────────────
\definecolor{demo@darkbg}{HTML}{1A1A2E}
\definecolor{demo@cyan}{HTML}{00D4AA}
\definecolor{demo@amber}{HTML}{F5A623}
\definecolor{demo@red}{HTML}{E74C3C}
\definecolor{demo@blue}{HTML}{3498DB}
\definecolor{demo@purple}{HTML}{9B59B6}
\definecolor{demo@green}{HTML}{2ECC71}

% ── tcolorbox styles ─────────────────────────────────────────────────────────
\newtcolorbox{toolbox}[1]{
    colback=sintefgrey,
    colframe=maincolor,
    fonttitle=\bfseries\small,
    title={#1},
    boxrule=0.5pt,
    arc=2pt,
    left=4pt, right=4pt, top=2pt, bottom=2pt,
}

\newtcolorbox{scenariobox}[2]{
    colback=sintefgrey,
    colframe=#1,
    fonttitle=\bfseries\small,
    title={#2},
    boxrule=0.8pt,
    arc=3pt,
    left=4pt, right=4pt, top=2pt, bottom=2pt,
}

\newtcolorbox{highlightbox}[1]{
    colback=white,
    colframe=#1,
    boxrule=1pt,
    arc=2pt,
    left=4pt, right=4pt, top=2pt, bottom=2pt,
}

% ── Code listing style ────────────────────────────────────────────────────────
\lstset{
    basicstyle=\ttfamily\scriptsize,
    keywordstyle=\color{demo@blue},
    commentstyle=\color{gray},
    stringstyle=\color{demo@red},
    showstringspaces=false,
    breaklines=true,
    frame=single,
    tabsize=4,
    language=bash
}

% ── Star rating command ───────────────────────────────────────────────────────
\newcommand{\stars}[1]{%
    \foreach \i in {1,...,#1}{\faStar}%
    \ifnum#1<5\relax
        \foreach \i in {1,...,\numexpr5-#1\relax}{\textcolor{gray!30}{\faStar}}%
    \fi
}

% ── Presentation metadata ────────────────────────────────────────────────────
\title{Politecnico di Torino Beamer Presentation Theme}
\subtitle{Using \LaTeX\ to prepare slides}
\course{Master's Degree in Computer Science}
\author{\href{mailto:author@studenti.polito.it}{Author Name}}
\IDnumber{1234567}
\email{author@studenti.polito.it}
\date{\today}

% ══════════════════════════════════════════════════════════════════════════════
\begin{document}
\maketitle

% ── Sections ──────────────────────────────────────────────────────────────────
% 01-introduction.tex - Introduction & Credits

\section{Introduction}

\begin{frame}{Beamer vs. PowerPoint}
Compared to PowerPoint, using \LaTeX\ with Beamer is better because:
\begin{itemize}
    \item \textbf{What-You-Mean-Is-What-You-Get}: You write the content, the computer handles the professional typesetting.
    \item \textbf{PDF Portability}: Produces a \texttt{.pdf} directly --- no worries about missing fonts, broken formulas, or mismatched software versions.
    \item \textbf{Consistency}: Extremely easy to maintain consistent layouts, fonts, highlights, and colors.
    \item \textbf{Mathematical Typesetting}: Standard \TeX\ math rendering is the best in class:
    \begin{equation*}
        \mathrm{i}\,\hslash\frac{\partial}{\partial t} \Psi(\mathbf{r},t) =
        -\frac{\hslash^2}{2\,m}\nabla^2\Psi(\mathbf{r},t) + V(\mathbf{r})\Psi(\mathbf{r},t)
    \end{equation*}
\end{itemize}
\end{frame}

\begin{frame}{Template Credits}
This presentation theme has a rich history of contributions:
\begin{itemize}
    \item \textbf{Original SINTEF Theme}: Created by \hrefcol{mailto:federico.zenith@sintef.no}{Federico Zenith}, derived from Håvard Berland's \texttt{beamerthementnu}.
    \item \textbf{PoliTo \& Sapienza Adaptation}: Adapted for Politecnico di Torino by \hrefcol{mailto:mattia.ippoliti@studenti.polito.it}{Mattia Ippoliti} and Sapienza by Andrea Gasparini.
    \item \textbf{College Beamer inspiration}: Liu Qilong's collection is retained as a documented influence; no current file is vendored from that repository.
    \item \textbf{Apograph Library Generalization}: Packaged and generalized by Elia Innocenti (providing robust fallbacks, path cleaning, and dynamic metadata).
\end{itemize}
\vspace{0.3cm}
\begin{center}
    \footnotesize Presentation content carries CC BY 4.0 attribution; the theme code is GPL-3.0-or-later. See \texttt{NOTICE} and \texttt{LICENSES/}.
\end{center}
\end{frame}

% 02-getting-started.tex - Preamble & Metadata Setup

\section{Getting Started}

\begin{frame}[fragile]{Getting Started}
\framesubtitle{Minimum PoliTo Beamer Document}
To start working with this theme, initialize your \LaTeX\ document with:
\begin{block}{Minimum Document Structure}
\begin{lstlisting}[language=TeX]
\documentclass[aspectratio=169]{beamer}
\usetheme{apographpolito}
% config.tex — the primary customization surface

\title{Presentation Title}
\subtitle{A concise subtitle}
\author{Author Name}
\date{\today}

\apographCourse{Course, event, or organization}
\apographStudentID{}
\apographEmail{author@example.com}
\apographSetClosingMessage{Thank you for listening!}

% Optional user-provided background. Uncomment to enable it; the theme still
% compiles with a fallback when the file is absent.
% \apographTitleBackground*{theme/assets/background}

% Use \apographThemeColor{main} for a dark document-wide color mode.
\apographThemeColor{white}

\begin{document}
  \begin{frame}{Hello, world!}
  \end{frame} %
\end{document}
\end{lstlisting}
\end{block}
\end{frame}

\begin{frame}[fragile]{Preamble Metadata}
\framesubtitle{Title Page Configuration}
Define presentation details in your preamble using the following commands:
\begin{block}{Metadata Commands}
\begin{lstlisting}[language=TeX]
\title{Sample Presentation Title}
\subtitle{Optional Subtitle}
\apographCourse{Master's Degree in Computer Science}
\author{Author Name}
\apographStudentID{1234567}
\apographEmail{author@studenti.polito.it}
\date{\today}
\end{lstlisting}
\end{block}
Render the title slide using \verb|\maketitle|.
\end{frame}

\begin{frame}[fragile]{Title Background Layouts}
\framesubtitle{Full Background vs. Split View}
You can style the title slide background using \verb|\apographTitleBackground| before \verb|\maketitle|:
\begin{itemize}
    \item \textbf{Full-Screen Background}:
    \begin{block}{Full Background Code}
    \verb|\apographTitleBackground{theme/assets/background}|
    \end{block}
    \item \textbf{Split Background (Right-Clipped)}:
    \begin{block}{Starred Split Background Code}
    \verb|\apographTitleBackground*{theme/assets/background}|
    \end{block}
\end{itemize}
\end{frame}

% 03-colors-styles.tex - Theme Slide Colors & footlinecolor

\section{Colors \& Styles}

\begin{frame}[fragile]{Theme Color Modes}
\framesubtitle{Preamble Configuration}
You can set the global slide background and text colors in the preamble using \verb|\themecolor{style}|:
\begin{itemize}
    \item \textbf{Light Theme (Default)}:
    \begin{block}{Light Theme Code}
    \verb|\themecolor{white}|
    \end{block}
    Sets a white background with dark grey text and maincolor accents.
    \item \textbf{Dark / Negative Theme}:
    \begin{block}{Dark Theme Code}
    \verb|\themecolor{main}|
    \end{block}
    Inverts the style: white text on a maincolor (PoliTo blue) background.
\end{itemize}
\end{frame}

\footlinecolor{sintefyellow}
\begin{frame}[fragile]{Footline Colors}
\framesubtitle{Customizing Slide Footers}
You can change the footline color dynamically by placing \verb|\footlinecolor{<color>}| before a frame:
\begin{itemize}
    \item Setting a color adds a footer block displaying the page number and presentation payoff (author | title).
    \item Current slide's footer is colored using \testcolor{sintefyellow}.
    \item Restore the default transparent footline with \verb|\footlinecolor{}| before a new frame.
\end{itemize}
\begin{block}{Usage}
\verb|\footlinecolor{sintefyellow}| \\
\verb|\begin{frame}{Footline Colors}...|
\end{block}
\end{frame}

\footlinecolor{}
\begin{frame}[fragile]{Built-in Colors}
\framesubtitle{Color Registry from sintefcolor.sty}
The package defines a set of official, harmonized colors for blocks and text:
\begin{itemize}
    \item \textbf{Primary Colors}:
    \testcolor{maincolor} and its sidebar companion \testcolor{sintefgrey} (or \texttt{sintefgray}).
    \item \textbf{Greens}:
    \testcolor{sinteflightgreen}, \testcolor{sintefgreen}, \testcolor{sintefdarkgreen}.
    \item \textbf{Accent Colors}:
    \testcolor{sintefyellow}, \testcolor{sintefred}, \testcolor{sinteflilla}.
    \item \textbf{Deprecated Colors} (for backward compatibility):
    \testcolor{sintefcyan}, \testcolor{sintefmagenta}, \testcolor{sinteflightgrey}.
\end{itemize}
\end{frame}

% 04-layout-elements.tex - Lists, Columns, Blocks, Fonts

\section{Layout Elements}

\begin{frame}[fragile]{Lists \& Columns}
\framesubtitle{Sequential Uncovering \& Grids}
\begin{columns}[T]
    \begin{column}{0.48\textwidth}
        \textbf{Sequential Lists}
        \begin{itemize}[<+->]
            \item Lists can reveal items step-by-step.
            \item Use the \verb|[<+->]| overlay option.
            \item Helps the audience focus on the current topic.
        \end{itemize}
    \end{column}
    \begin{column}{0.48\textwidth}
        \textbf{Two-Column Splits}
        \begin{block}{Multi-column layout code}
\begin{lstlisting}[language=TeX]
\begin{columns}
  \begin{column}{0.5\textwidth}
    Left content...
  \end{column}
  \begin{column}{0.5\textwidth}
    Right content...
  \end{column}
\end{columns}
\end{lstlisting}
        \end{block}
    \end{column}
\end{columns}
\end{frame}

\begin{frame}[fragile]{Standard vs. Colour Blocks}
\framesubtitle{Organizing Key Points}
\begin{columns}[T]
    \begin{column}{0.48\textwidth}
        \begin{block}{Standard Block}
        Standard blocks adjust color automatically based on the footline color (falls back to grey/blue).
\begin{lstlisting}[language=TeX]
\begin{block}{Title}
Content...
\end{block}
\end{lstlisting}
        \end{block}
    \end{column}
    \begin{column}{0.48\textwidth}
        \begin{colorblock}[black]{sinteflightgreen}{Colour Block}
        You pick the color explicitly. Text is white by default, but customizable via the optional argument.
\begin{lstlisting}[language=TeX]
\begin{colorblock}[black]
  {sinteflightgreen}
  {Title}
  Content...
\end{colorblock}
\end{lstlisting}
        \end{colorblock}
    \end{column}
\end{columns}
\end{frame}

\begin{frame}{Typography \& Fonts}
\framesubtitle{Maximizing Readability}
The template defaults to Carlito (Calibri clone) sans-serif fonts:
\begin{itemize}
    \item \textbf{Sans-serif vs. Serif}:
    \begin{itemize}
        \item \textsf{Sans-serif fonts} are highly recommended for general screens and projectors.
        \item \textrm{Use serif fonts only for formulas, mathematics, or extremely high-definition displays}.
    \end{itemize}
    \item \textbf{Style Warnings}:
    \begin{itemize}
        \item \texttt{Never use monospace fonts for paragraphs or normal body text.} Keep it for terminal commands or code blocks.
        \item {\frakfamily Gothic, calligraphic, or decorative fonts should be completely avoided.}
    \end{itemize}
\end{itemize}
\end{frame}

% 05-custom-boxes.tex - Custom Boxes & Macros

\section{Custom Containers}

\begin{frame}[fragile]{Custom tcolorbox Containers}
\framesubtitle{Custom styled box components}
This package defines three specialized container environments for code, notes, and highlights:
\begin{columns}[T]
    \begin{column}{0.48\textwidth}
        \begin{toolbox}{\faBriefcase\; Toolbox Box}
            \footnotesize
            Created via \texttt{\textbackslash begin\{toolbox\}\{Title\}}.
            Uses \texttt{maincolor} frame and \texttt{sintefgrey} background.
        \end{toolbox}
        
        \vspace{0.2cm}
        
        \begin{highlightbox}{demo@purple}
            \footnotesize
            \textbf{\faInfoCircle\; Highlight Box}
            
            Created via \texttt{\textbackslash begin\{highlightbox\}\{color\}}.
            Uses white background and a customizable frame color.
        \end{highlightbox}
    \end{column}
    \begin{column}{0.48\textwidth}
        \begin{scenariobox}{demo@amber}{\faExclamationTriangle\; Scenario Box}
            \footnotesize
            Created via \texttt{\textbackslash begin\{scenariobox\}\{color\}\{Title\}}.
            Uses custom frame color, custom title, and \texttt{sintefgrey} background.
        \end{scenariobox}
        
        \vspace{0.1cm}
        \scriptsize
        These containers support icons (e.g. \texttt{fontawesome5}) and nested lists to summarize key items visually.
    \end{column}
\end{columns}
\end{frame}

\begin{frame}[fragile]{Rating Stars Command}
\framesubtitle{Difficulty or Importance rating macro}
You can use the custom \verb|\stars{N}| command to display a star rating out of 5:
\begin{itemize}
    \item \textbf{1 Star}: \stars{1}
    \item \textbf{3 Stars}: \stars{3}
    \item \textbf{5 Stars}: \stars{5}
\end{itemize}

\begin{block}{Example Code}
\begin{lstlisting}[language=TeX]
\begin{frame}{Advanced Slide \hfill \stars{4}}
  This slide is of advanced difficulty level.
\end{frame} %
\end{lstlisting}
\end{block}
\end{frame}

% 06-diagrams-tables.tex - Diagrams & Tables

\section{Diagrams \& Tables}

\begin{frame}{Drawing Flowcharts via TikZ}
\framesubtitle{Native vector flows using TikZ library}
\centering
\begin{tikzpicture}[
    nodeDistance/.style={shape=rectangle, rounded corners=4pt, draw=maincolor, fill=sintefgrey, minimum width=2.4cm, minimum height=1.0cm, font=\scriptsize\bfseries, align=center},
    arrow/.style={-{Stealth[length=3mm]}, thick, maincolor}
]
    % Nodes
    \node[nodeDistance] (step1) {\faFile*[regular]\\[0.1cm] 1. Draft Content};
    \node[nodeDistance, right=1.2cm of step1] (step2) {\faCode\\[0.1cm] 2. Typeset};
    \node[nodeDistance, right=1.2cm of step2] (step3) {\faEye\\[0.1cm] 3. Preview PDF};

    % Arrows
    \draw[arrow] (step1) -- (step2);
    \draw[arrow] (step2) -- (step3);
\end{tikzpicture}

\vspace{0.4cm}
\begin{columns}[T]
    \begin{column}{0.48\textwidth}
        \footnotesize
        \textbf{Workflow Design}:
        \begin{itemize}
            \item Position nodes dynamically with coordinates.
            \item Style frames using the theme colors (e.g. \texttt{maincolor}, \texttt{sintefgrey}).
        \end{itemize}
    \end{column}
    \begin{column}{0.48\textwidth}
        \footnotesize
        \textbf{Vector Quality}:
        \begin{itemize}
            \item Native TikZ lines scale losslessly to any resolution.
            \item Avoids importing fuzzy screenshot images.
        \end{itemize}
    \end{column}
\end{columns}
\end{frame}

\begin{frame}{State Transitions \& Gating}
\framesubtitle{Showing before/after state flows}
\begin{columns}[T]
    \begin{column}{0.45\textwidth}
        \centering
        \vspace{0.2cm}
        \begin{tikzpicture}[
            box/.style={draw=maincolor, fill=sintefgrey, rounded corners=3pt,
                minimum width=4.2cm, minimum height=0.7cm,
                font=\scriptsize, align=center, text=darkgray},
            locked/.style={draw=demo@red, fill=demo@red!8, rounded corners=3pt,
                minimum width=4.2cm, minimum height=0.7cm,
                font=\scriptsize, align=center, text=demo@red},
            unlocked/.style={draw=demo@green, fill=demo@green!8, rounded corners=3pt,
                minimum width=4.2cm, minimum height=0.7cm,
                font=\scriptsize, align=center, text=demo@green!70!black},
            arrow/.style={-{Stealth[length=2.5mm]}, thick, maincolor},
        ]
            \node[locked] (lock) {\faLock\; Panel State --- \textbf{LOCKED}};
            \node[box, below=0.6cm of lock] (cmd) {\faTerminal\; \texttt{run\_command\_here}};
            \node[unlocked, below=0.6cm of cmd] (unlock) {\faLockOpen\; Panel State --- \textbf{UNLOCKED}};

            \draw[arrow] (lock) -- node[right, font=\tiny, text=gray] {action required} (cmd);
            \draw[arrow, demo@green] (cmd) -- node[right, font=\tiny, text=demo@green!70!black] {unlock trigger} (unlock);
        \end{tikzpicture}
    \end{column}
    \begin{column}{0.50\textwidth}
        \begin{highlightbox}{maincolor}
            \footnotesize\bfseries
            \faInfoCircle\; Gating Mechanics
            \vspace{0.15cm}

            \normalfont\footnotesize
            You can design state transition workflows using colors:
            \begin{itemize}
                \item Use \testcolor{demo@red} for restricted or locked states.
                \item Use \testcolor{demo@green} for active or resolved states.
                \item Link processes step-by-step using TikZ paths.
            \end{itemize}
        \end{highlightbox}
    \end{column}
\end{columns}
\end{frame}

\begin{frame}{Tables with Icons}
\framesubtitle{Tabular layout demonstrations}
\begin{toolbox}{\faTable\; Structured Information Table}
    \footnotesize
    \renewcommand{\arraystretch}{1.3}
    \begin{tabularx}{\textwidth}{@{}c l X@{}}
        \toprule
        \textbf{Icon} & \textbf{Component} & \textbf{Functionality} \\
        \midrule
        \faDesktop & GUI Workspace & Visualizes outputs, file systems, and details. \\
        \faTerminal & CLI Terminal & Interactive low-level command workspace execution. \\
        \faBookOpen & Notebook & Tracks objectives, logs findings, and saves state progress. \\
        \bottomrule
    \end{tabularx}
\end{toolbox}
\end{frame}

% 07-special-slides.tex - Special Slides & Backmatter

\section{Special Slides}

\begin{apographchapter}[theme/assets/background_negative]{}{Chapter Divider Slides}
\begin{itemize}
    \item Chapter slides split the slide layout: the right side contains a clipped image, while the left side displays section headers.
    \item Useful to structure transitions between long sections.
\end{itemize}
\end{apographchapter}

\begin{apographsidepic}{theme/assets/background_alternative}{Side-Picture Slides}
\begin{itemize}
    \item Use the \texttt{apographsidepic} environment with image and title arguments.
    \item Splits the screen horizontally: left side contains text content (takes up to 60\% width), right side contains a full-height image.
\end{itemize}
\end{apographsidepic}

\begin{frame}[fragile]{Ending the Presentation}
\framesubtitle{Backmatter payoffs and acknowledgments}
To insert a final payoff/thank-you slide:
\begin{itemize}
    \item Call \verb|\apographBackmatter| at the very end of your document:
    \begin{block}{Default Backmatter}
    \verb|\apographBackmatter|
    \end{block}
    This displays the presentation title and standard thank-you text with author contact email.
    \item To hide the title on the final slide, pass the optional \verb|notitle| parameter:
    \begin{block}{No-Title Backmatter}
    \verb|\apographBackmatter[notitle]|
    \end{block}
\end{itemize}
\end{frame}


\title{} % Reset for backmatter
\backmatter

\end{document}
